\section{Metodologia documental e organização do repositório}
Na revisão documental de 8 de novembro de 2026, cada uma das 66 indicações recebidas ganhou identificador estável e parecer individual, publicados no repositório do autor \cite{souzaRepositorio2026}. A matriz distingue autenticação de representação do riso, referência lexical ou temática, identificação bibliográfica, recuperação de exemplar específico e casos sem confirmação. As buscas literais das 22 frases e suas limitações constam em arquivos próprios. A existência de um catálogo não foi tratada como substituto de passagem primária; a pesquisa de OCR não foi descrita como leitura visual integral. O método estatístico a seguir permanece prospectivo: as novas aquisições não transformam a coleta dirigida em amostra probabilística.
\subsection{Desenho executado, escopo e material de partida}
O estudo executado é exploratório, documental e crítico, com busca dirigida e coleta não probabilística. Seu objeto empírico tem dois níveis: as alegações e supostas citações presentes no material conversacional fornecido pelo autor; e os documentos recuperados para verificá-las ou contextualizá-las. Não é uma revisão sistemática exaustiva, nem uma amostra representativa de toda a escrita portuguesa entre 1872 e 2026. O recorte de 1900 a 2009 organiza a cobertura anual do núcleo histórico sugerido no material; testemunhos anteriores e posteriores funcionam como extensões documentadas.

A consulta foi realizada em 8 de novembro de 2026. Utilizaram-se buscas por autor, título, variante gráfica e ano em acervos institucionais e bibliotecas digitais, incluindo Project Gutenberg, Wikisource, o corpus de Machado de Assis mantido na Brigham Young University, a Hemeroteca Digital Brasileira e repositórios universitários. Variantes sem acento, pontuação alterada e repetições distintas foram consideradas como pistas. Resultados de buscadores serviram para descoberta; a passagem recuperada, a imagem ou o texto integral acessível fundamentaram a classificação de evidência.

O material conversacional não foi contado como ocorrência histórica. Ele é um roteiro de investigação que contém hipóteses, exemplos, contradições e transcrições ainda sem autenticação. A auditoria conserva o assunto de cada trecho, distinguindo a formulação original de sua avaliação científica. Os arquivos integrais de trabalho permanecem locais; a versão pública contém uma matriz de alegações com os elementos necessários à conferência, sem divulgar a conversa inteira ou metadados pessoais desnecessários.

\subsection{Unidades e metadados}
Distinguem-se quatro unidades. A \emph{referência bibliográfica} identifica uma obra; o \emph{documento amostral} é uma edição, número de periódico ou sessão de interação; a \emph{ocorrência} é uma passagem contextualizada; o \emph{tipo gráfico} é uma categoria, como \forma{quiá} ou \forma{kkkk}. Três passagens do mesmo romance não equivalem a três documentos independentes. Uma dissertação que reproduz um chat é o documento consultado, enquanto o chat é o testemunho reportado. A unidade não deve mudar silenciosamente entre tabelas.

Para cada testemunho registram-se identificador, autor ou periódico, título, ano da obra, ano da edição ou estudo, data da interação quando disponível, país, suporte, grafia literal, contexto, localizador, URL, data de acesso, natureza da evidência e restrições. Em PDFs, informa-se a página do arquivo e, quando conhecida, a paginação impressa. Em textos eletrônicos sem paginação estável, utiliza-se capítulo e âncora textual. Um identificador baseado em sequência exata permite reencontrar a passagem sem inventar uma página.

A grafia diplomática conserva acento, caixa e pontuação na medida permitida pelo testemunho. A normalização é um campo separado, usado para localizar variantes. ``Quá quá quá'' e ``quá-quá-quá'' podem ser agrupados numa análise de tipos, mas não são apresentados como transcrições idênticas. Um OCR que lê uma vogal incorretamente não decide a forma final: a imagem prevalece quando disponível.

\subsection{Níveis de evidência e regras de inclusão}
Adotam-se os códigos E1, fonte primária com imagem inspecionada e localizador; E2, texto primário em edição ou transcrição eletrônica consultada, sem colação integral com o impresso; E3, interação reproduzida em estudo acadêmico, com página e contexto identificados; E4, alegação secundária, paráfrase ou referência sem passagem autenticada. E1 não significa ausência de qualquer erro editorial; E2 não equivale a edição crítica. Esses códigos qualificam a cadeia de acesso, não a importância literária.

Incluem-se no inventário de atestações passagens cujo contexto sustenta leitura de riso e cujo suporte pode ser identificado. Excluem-se da contagem principal nomes comerciais, títulos de programas, siglas e sequências de letras sem função de riso demonstrada. A palavra ``risos'' numa descrição não é automaticamente rubrica de transcrição. Uma cadeia \forma{KKKK} que nomeia um empreendimento ou uma organização não é uma gargalhada digital. Tais falsos positivos são registrados quando relevantes para explicar a triagem.

As frases literais do material inicial não foram autenticadas por existirem livros com os mesmos títulos. A identificação da obra e a verificação da passagem são etapas distintas. A ausência de correspondência numa busca eletrônica também não prova falsificação: pode haver diferenças de edição, OCR ou redação. O estatuto adequado é ``não autenticada nesta coleta'', salvo quando há uma contradição documental específica.

\subsection{Rastreabilidade, downloads e disponibilidade}
O repositório do autor, \repo, concentra o manuscrito, bibliografia, dados e scripts \cite{souzaRepositorio2026}. O inventário de aquisição registra URL, resultado de acesso, tamanho e SHA-256 dos arquivos efetivamente baixados. Falhas de acesso são mantidas no histórico. Os scripts distinguem arquivos obtidos localmente de documentos consultados por serviço de leitura web; uma consulta bem-sucedida num ambiente não é apresentada como download bem-sucedido em outro.

O GitHub é infraestrutura de versionamento e disponibilização, não prova de validade da hipótese. A referência metodológica identifica o autor e o repositório; a reprodução deve usar o commit ou a versão efetivamente consultada. Links, checksums e scripts contribuem para localização e reuso segundo os princípios FAIR \cite{wilkinson2016}. Não se presume que acesso gratuito autorize republicação integral. Documentos com direitos não esclarecidos permanecem locais, com metadados públicos e instruções de obtenção. Arquivos redistribuídos conservam fonte e condições aplicáveis.

\subsection{Controle de qualidade e limites de validação}
Efetuou-se conferência contextual das passagens selecionadas, inspeção visual da capa de \textit{O Tico-Tico} e validação dos localizadores disponíveis. Uma revisão eletrônica verifica chaves bibliográficas, referências cruzadas, ausência de citações não resolvidas, contagem de registros e consistência entre tabelas geradas e arquivos de dados. A revisão é feita em uma etapa de trabalho assistida por ferramentas; não se declara dupla codificação independente ou concordância entre avaliadores que não ocorreu.

Uma ampliação deverá empregar pelo menos dois anotadores em uma subamostra, medir acordo por categoria e discutir divergências. Textos antigos exigem atenção a acentos, hifens, fontes tipográficas e ordem das colunas. A validação humana não deve limitar-se a confirmar a regex: precisa verificar a função da sequência e sua atribuição ao documento correto.

\subsection{Síntese quantitativa, inventário de bases e links de reprodução}
O quadro quantitativo reúne \CorpusSize{} testemunhos, 66 indicações auditadas e 22 citações alegadas verificadas. As 57 referências da seleção por acervo e década e os 17 títulos inéditos incluem documentos que também integram o núcleo. As tabelas indicam as unidades e os arquivos correspondentes no GitHub.

Os dezesseis endereços indicados pelo autor não equivalem a dezesseis bases independentes: o APESP aparece como portal e item individual; a BBM, como coleção específica e acervo geral; e o Orkut é conteúdo preservado na própria Wayback. O inventário conserva todos os endereços e explicita essas relações. Os recursos complementares, como bibliotecas digitais e repositórios universitários, estão identificados por URL individual no manifesto, com inventário adicional de seus domínios. Um domínio técnico, uma página de catálogo e uma obra consultada não são unidades intercambiáveis.

O manifesto histórico inclui arquivos auxiliares, metadados, código de interfaces, documentos sem ocorrência autenticada e o histórico de uma referência retirada. As quantidades de aquisições e falhas constam no quadro quantitativo. A unidade do manifesto é a entrada de aquisição.

\small\begin{longtable}{>{\raggedright\arraybackslash}p{.31\textwidth}>{\raggedright\arraybackslash}p{.13\textwidth}>{\raggedright\arraybackslash}p{.45\textwidth}}
\caption{Balanço quantitativo consolidado do levantamento. As linhas têm unidades distintas e não devem ser somadas.}\\
\toprule Indicador & Quantidade & Unidade, escopo e dados\\\midrule\endfirsthead
\toprule Indicador & Quantidade & Unidade, escopo e dados\\\midrule\endhead
Endereços indicados pelo autor & 16 & endereços. Portais, coleções, item individual e captura; há sobreposições institucionais. \href{https://github.com/murilloed/historia-riso-escrito-portugues/blob/main/data/base_registry.csv}{Dados no GitHub} \\
Núcleo de representações autenticadas & 121 & documentos/testemunhos. Inclui documentos primários e interações reportadas; exclui menções apenas lexicais. \href{https://github.com/murilloed/historia-riso-escrito-portugues/blob/main/data/attestations.csv}{Dados no GitHub} \\
Níveis do núcleo E1 / E2 / E3 & 23 / 95 / 3 & testemunhos. Imagem / texto primário eletrônico / interação reportada; partição dos 121 registros. \href{https://github.com/murilloed/historia-riso-escrito-portugues/blob/main/data/attestations.csv}{Dados no GitHub} \\
Indicações iniciais com parecer & 66 & indicações. Auditoria bibliográfica; não são 66 ocorrências positivas. \href{https://github.com/murilloed/historia-riso-escrito-portugues/blob/main/data/references_audit.csv}{Dados no GitHub} \\
Citações alegadas submetidas a verificação & 22 & alegações. Nenhuma autenticada integralmente; os resultados constam no parecer. \href{https://github.com/murilloed/historia-riso-escrito-portugues/blob/main/data/claims_verification.csv}{Dados no GitHub} \\
Referências por acervo e década & 135 & referências selecionadas. Inclui menções lexicais, rubricas e grafias miméticas, com fontes alternativas identificadas na matriz. \href{https://github.com/murilloed/historia-riso-escrito-portugues/blob/main/data/archive_decade_references.csv}{Dados no GitHub} \\
Linhas de acompanhamento por acervo/década & 137 & linhas. 135 linhas com referência e dois itens conferidos sem ocorrência; não é o universo de documentos. \href{https://github.com/murilloed/historia-riso-escrito-portugues/blob/main/data/archive_decade_matrix.csv}{Dados no GitHub} \\
Títulos inéditos selecionados & 17 & títulos. Rodada 4: 14 Brasiliana e três espaços digitais; duas fontes já incorporadas ao núcleo. \href{https://github.com/murilloed/historia-riso-escrito-portugues/blob/main/data/revision4_new_titles.csv}{Dados no GitHub} \\
Anos no inventário anual do núcleo & 155 & anos. 1872-2026; anos sem testemunho não significam ausência histórica. \href{https://github.com/murilloed/historia-riso-escrito-portugues/blob/main/data/annual_coverage.csv}{Dados no GitHub} \\
Entradas no manifesto histórico & 1635 & entradas de aquisição. Inclui apoio, metadados, recursos técnicos, falhas e referências retiradas; não documentos lidos. \href{https://github.com/murilloed/historia-riso-escrito-portugues/blob/main/data/downloads.json}{Dados no GitHub} \\
Entradas com conteúdo baixado & 1482 & entradas não vazias. Inclui HTML, PDF, imagem, JSON e código de interface; não equivale a leitura integral. \href{https://github.com/murilloed/historia-riso-escrito-portugues/blob/main/data/downloads.json}{Dados no GitHub} \\
Respostas vazias preservadas no manifesto & 2 & entradas vazias. Não contadas como documentos analisados ou resultados negativos. \href{https://github.com/murilloed/historia-riso-escrito-portugues/blob/main/data/downloads.json}{Dados no GitHub} \\
Entradas com falha de aquisição & 151 & entradas com falha. Estado registrado por identificador; não representa todas as tentativas históricas. \href{https://github.com/murilloed/historia-riso-escrito-portugues/blob/main/data/downloads.json}{Dados no GitHub} \\
Textos integrais com cópia pública & 8 & cópias públicas. Demais arquivos integrais ficam locais quando a redistribuição não foi verificada. \href{https://github.com/murilloed/historia-riso-escrito-portugues/blob/main/sources/README.md}{Dados no GitHub} \\
Referências bibliográficas & 322 & entradas BibTeX. Inclui teoria, documentos, catálogos e portais; não corresponde ao número de atestações. \href{https://github.com/murilloed/historia-riso-escrito-portugues/blob/main/artigo/references.bib}{Dados no GitHub} \\
Cobertura anual de 1900 a 2009 & 110 & anos com atestação. Intervalo de 110 anos. A busca dirigida incluiu rubricas e grafias miméticas. \href{https://github.com/murilloed/historia-riso-escrito-portugues/blob/main/data/annual_coverage.csv}{Dados no GitHub} \\
Conjuntos de recursos de acesso indicados & 13 & grupos operacionais. 16 endereços agrupados: APESP portal/item, BBM coleção/acervo geral e Wayback/Orkut; não instituições independentes. \href{https://github.com/murilloed/historia-riso-escrito-portugues/blob/main/data/base_registry.csv}{Dados no GitHub} \\
Fontes da seleção por acervo/década & 129 & identificadores de fonte. As 135 seleções reutilizam documentos entre quotas; 129 source\_id distintos. \href{https://github.com/murilloed/historia-riso-escrito-portugues/blob/main/data/evidence_register.csv}{Dados no GitHub} \\
Fontes selecionadas nos quatro inventários & 258 & identificadores de fonte. Reconciliação por source\_id e URL; não independência de obras, edições ou testemunhos. \href{https://github.com/murilloed/historia-riso-escrito-portugues/blob/main/data/consolidated_inventory.json}{Dados no GitHub} \\
Rubricas / marcadores miméticos e digitais do núcleo & 96 / 25 & registros. Partição do núcleo; rubricas descrevem reações, sem transcrever necessariamente o som. \href{https://github.com/murilloed/historia-riso-escrito-portugues/blob/main/data/consolidated_inventory.json}{Dados no GitHub} \\
Alternativas literárias de período & 2 & fontes. Bocage 1800 e Gonzaga 1812; não ocorrências nos exemplares fixos da Gazeta. \href{https://github.com/murilloed/historia-riso-escrito-portugues/blob/main/data/period_alternatives.csv}{Dados no GitHub} \\
Itens conferidos sem ocorrência localizada & 2 & itens. Leitura concluída; não prova de ausência na série nem quota positiva. \href{https://github.com/murilloed/historia-riso-escrito-portugues/blob/main/data/archive_decade_matrix.csv}{Dados no GitHub} \\
Confirmação no acervo indicado / fonte alternativa & 65 / 70 & linhas de quota. Partição das 135 seleções; confirmação alternativa não equivale a cotejo no portal original. \href{https://github.com/murilloed/historia-riso-escrito-portugues/blob/main/data/archive_decade_matrix.csv}{Dados no GitHub} \\
Classes lexicais / rubricas / miméticas da seleção & 112 / 18 / 5 & seleções. Partição das 135 linhas; não prevalência histórica nem fontes distintas. \href{https://github.com/murilloed/historia-riso-escrito-portugues/blob/main/data/consolidated_inventory.json}{Dados no GitHub} \\
Domínios registrados no manifesto & 87 & hosts. Inclui subdomínios, serviços técnicos e aquisições falhas; não equivale a bases consultadas integralmente. \href{https://github.com/murilloed/historia-riso-escrito-portugues/blob/main/data/source_domains.csv}{Dados no GitHub} \\
Fontes anteriores a 1900 & 40 & identificadores de fonte. 47 vínculos nos quatro inventários; inclui documentos desde 1645. Não é cobertura anual contínua. \href{https://github.com/murilloed/historia-riso-escrito-portugues/blob/main/data/consolidated_inventory.json}{Dados no GitHub} \\
\bottomrule\end{longtable}\normalsize

\small\begin{longtable}{>{\raggedright\arraybackslash}p{.29\textwidth}>{\raggedright\arraybackslash}p{.38\textwidth}>{\raggedright\arraybackslash}p{.22\textwidth}}
\caption{Inventário consolidado dos 16 endereços indicados: portais, coleções, item e captura. Sobreposições não são novas bases independentes.}\\
\toprule ID / recurso & Relação entre os recursos & Seleção e rastreabilidade\\\midrule\endfirsthead
\toprule ID / recurso & Relação entre os recursos & Seleção e rastreabilidade\\\midrule\endhead
H01 DocVirt / Jardinópolis \cite{r3archiveH01} & Portal/coleção; ver limitações de acesso e cobertura. & 11; referências por acervo e década \href{https://app.docvirt.com/dvd_ajhpmj/pageid/1}{Fonte}; \href{https://github.com/murilloed/historia-riso-escrito-portugues/blob/main/data/archive_decade_references.csv}{GitHub} \\
H02 O Globo \cite{r3archiveH02} & Portal/coleção; ver limitações de acesso e cobertura. & 7; referências por acervo e década \href{https://oglobo.globo.com/acervo/}{Fonte}; \href{https://github.com/murilloed/historia-riso-escrito-portugues/blob/main/data/archive_decade_references.csv}{GitHub} \\
H03 Folha \cite{r3archiveH03} & Portal/coleção; ver limitações de acesso e cobertura. & 11; referências por acervo e década \href{https://acervo.folha.com.br/index.do}{Fonte}; \href{https://github.com/murilloed/historia-riso-escrito-portugues/blob/main/data/archive_decade_references.csv}{GitHub} \\
H04 BORE / USP, Jornais \cite{r3archiveH04} & Portal/coleção; ver limitações de acesso e cobertura. & 13; referências por acervo e década \href{https://obrasraras.usp.br/xmlui/handle/123456789/4262}{Fonte}; \href{https://github.com/murilloed/historia-riso-escrito-portugues/blob/main/data/archive_decade_references.csv}{GitHub} \\
H05 Hemeroteca / BN \cite{r3archiveH05} & Portal/coleção; ver limitações de acesso e cobertura. & 21; referências por acervo e década \href{https://memoria.bn.gov.br/hdb/periodico.aspx}{Fonte}; \href{https://github.com/murilloed/historia-riso-escrito-portugues/blob/main/data/archive_decade_references.csv}{GitHub} \\
H06 Correio Braziliense \cite{r3archiveH06} & Portal/coleção; ver limitações de acesso e cobertura. & 6; referências por acervo e década \href{https://www.correiobraziliense.com.br/acervo}{Fonte}; \href{https://github.com/murilloed/historia-riso-escrito-portugues/blob/main/data/archive_decade_references.csv}{GitHub} \\
H07 Google News / Jornal do Brasil \cite{r3archiveH07} & Portal/coleção; ver limitações de acesso e cobertura. & 7; referências por acervo e década \href{https://news.google.com/newspapers?nid=0qX8s2k1IRwC}{Fonte}; \href{https://github.com/murilloed/historia-riso-escrito-portugues/blob/main/data/archive_decade_references.csv}{GitHub} \\
H08 APESP / portal \cite{r3archiveH08} & Portal APESP; H13 é um item deste mesmo acervo. & 0; referências por acervo e década \href{https://atom.arquivoestado.sp.gov.br/}{Fonte}; \href{https://github.com/murilloed/historia-riso-escrito-portugues/blob/main/data/archive_decade_references.csv}{GitHub} \\
H09 Senado / acervo de notícias \cite{r3archiveH09} & Portal/coleção; ver limitações de acesso e cobertura. & 3; referências por acervo e década \href{https://www12.senado.leg.br/noticias/acervo-historico}{Fonte}; \href{https://github.com/murilloed/historia-riso-escrito-portugues/blob/main/data/archive_decade_references.csv}{GitHub} \\
H10 Diario de Pernambuco / portal \cite{r3archiveH10} & Portal/coleção; ver limitações de acesso e cobertura. & 16; referências por acervo e década \href{https://www.diariodepernambuco.com.br/acervo}{Fonte}; \href{https://github.com/murilloed/historia-riso-escrito-portugues/blob/main/data/archive_decade_references.csv}{GitHub} \\
H11 CEPE / periódicos \cite{r3archiveH11} & Portal/coleção; ver limitações de acesso e cobertura. & 11; referências por acervo e década \href{https://acervo.cepe.com.br/}{Fonte}; \href{https://github.com/murilloed/historia-riso-escrito-portugues/blob/main/data/archive_decade_references.csv}{GitHub} \\
H12 BBM / Gazeta, coleção 8427 \cite{r3archiveH12} & Coleção Gazeta na BBM; H14 é o acervo geral da mesma biblioteca. & 0; referências por acervo e década \href{https://digital.bbm.usp.br/handle/bbm/8427}{Fonte}; \href{https://github.com/murilloed/historia-riso-escrito-portugues/blob/main/data/archive_decade_references.csv}{GitHub} \\
H13 APESP / item D000071 \cite{r3archiveH13} & Item individual do APESP; não é outra base independente. & 0; referências por acervo e década \href{https://atom.arquivoestado.sp.gov.br/br-spapesp-hemapesp-s001-b000565-d000071}{Fonte}; \href{https://github.com/murilloed/historia-riso-escrito-portugues/blob/main/data/archive_decade_references.csv}{GitHub} \\
H14 Brasiliana / BBM: acervo geral \cite{bbmAcervo2026} & Mesma biblioteca de H12; coleção geral, não outra instituição independente. & 14; títulos inéditos selecionados, rodada 4 \href{https://digital.bbm.usp.br/handle/bbm/1}{Fonte}; \href{https://github.com/murilloed/historia-riso-escrito-portugues/blob/main/data/revision4_new_titles.csv}{GitHub} \\
H15 Orkut preservado na Wayback \cite{r4orkutindex} & Conteúdo do Orkut preservado no serviço H16; não plataforma de preservação independente. & 1; comunidade selecionada, rodada 4 \href{https://web.archive.org/web/20141001005309/http://orkut.google.com/}{Fonte}; \href{https://github.com/murilloed/historia-riso-escrito-portugues/blob/main/data/revision4_new_titles.csv}{GitHub} \\
H16 Wayback Machine: portal geral \cite{waybackPortal2026} & Blog e fórum; Orkut contado separadamente em H15. O portal não foi pesquisado exaustivamente. & 2; blog e fórum selecionados, rodada 4 \href{https://web.archive.org/}{Fonte}; \href{https://github.com/murilloed/historia-riso-escrito-portugues/blob/main/data/revision4_new_titles.csv}{GitHub} \\
\bottomrule\end{longtable}\normalsize


Um zero na coluna de seleção indica ausência de referência autenticada naquele recorte e rodada, não ausência de consulta ao portal nem inexistência histórica de riso.

Os dados completos estão em \href{https://github.com/murilloed/historia-riso-escrito-portugues/blob/main/data/base_registry.csv}{bases}, \href{https://github.com/murilloed/historia-riso-escrito-portugues/blob/main/data/method_quantification.csv}{quantificação}, \href{https://github.com/murilloed/historia-riso-escrito-portugues/blob/main/data/source_domains.csv}{domínios} e \href{https://github.com/murilloed/historia-riso-escrito-portugues/blob/main/data/downloads.json}{aquisições}. Os links levam à versão corrente; a reprodução deve identificar a versão ou o commit consultado. O inventário explicita pendências e não equivale à cobertura integral das coleções.

\section{Fundamentação estatística da quantidade de documentos por ano}
\subsection{Por que três referências não constituem justificativa estatística}
O roteiro inicial pede três jornais e três livros por período. Esse número é uma quota editorial útil para organizar exemplos, mas não fornece representatividade por si só. A existência de três títulos conhecidos não informa o número de páginas examinadas, a quantidade de oportunidades de riso, o universo de documentos preservados ou a dependência entre passagens. A repetição de exemplos não estima frequência social. Tampouco um cálculo de amostra transforma retrospectivamente a busca dirigida em sorteio probabilístico.

É preciso separar dois objetivos. Para demonstrar \emph{existência}, uma atestação autêntica basta para aquele documento e data. Para estimar \emph{prevalência} ou comparar anos, é necessário definir população, unidade amostral, mecanismo de seleção e incerteza. Para demonstrar \emph{origem}, nenhum tamanho de amostra finito elimina automaticamente a possibilidade de documentos anteriores perdidos ou não inspecionados. A estatística auxilia o dimensionamento e a descrição; não produz uma certidão de nascimento etimológico.

\subsection{Desenho prospectivo estratificado por ano e suporte}
Propõe-se construir uma moldura de documentos para cada ano $y$ e suporte $s$, distinguindo ao menos livros e periódicos. O estrato é $h=(y,s)$ e contém $N_h$ unidades elegíveis. Livros podem ser definidos por edição publicada naquele ano; periódicos, por número, com páginas selecionadas em segundo estágio quando necessário. Interações digitais formam estratos próprios, pois sessões, mensagens e usuários não são equivalentes a volumes impressos. País, gênero e acervo devem ser registrados para controlar composição.

O estimando básico sugerido é a proporção de documentos do estrato com pelo menos uma representação gráfica de riso: $p_h=P(I_h=1)$. Outro estimando possível é a frequência de determinado tipo por quantidade de palavras inspecionadas. Esses objetos não se confundem. Uma obra longa tende a oferecer mais oportunidades de ocorrência, razão pela qual também se registra extensão. A frequência normalizada, se houver denominadores confiáveis, será $f_h=10^6 C_h/W_h$, em que $C_h$ é número de ocorrências e $W_h$ o total de palavras elegíveis.

Quando os estratos têm dimensões conhecidas, pode-se fazer censo dos pequenos e sorteio dos grandes, registrando probabilidades de inclusão. A estimativa agregada usa pesos proporcionais aos universos, e não uma média simples de estratos com tamanhos diferentes. Se custos e variâncias forem conhecidos, a alocação ótima é proporcional a $N_hS_h/\sqrt{c_h}$, em que $S_h$ é o desvio-padrão do indicador e $c_h$ o custo por unidade. Sem piloto ou moldura, essa alocação é uma proposta, não um procedimento executado \cite{cochran1977}.

\subsection{Dimensionamento para uma proporção}
Sob amostragem aleatória simples, população grande, nível nominal de 95\% e erro absoluto $e$, um dimensionamento inicial é
\begin{equation}
n_0=\frac{z_{0,975}^2p(1-p)}{e^2},\qquad z_{0,975}=1{,}96.
\end{equation}
Na ausência de proporção piloto, $p=0{,}5$ maximiza a variância. Com $e=0{,}10$, resulta $n_0=96{,}04$, arredondado para 97 documentos \emph{por estrato} se a precisão for exigida em cada combinação de ano e suporte. Para $e=0{,}05$, são 385. Esses números não significam 97 ou 385 referências bibliográficas a citar no artigo: as referências identificam os documentos examinados, enquanto tabelas exemplificativas podem exibir apenas uma seleção transparente.

Para universo finito conhecido, a correção é
\begin{equation}
n_h=\left\lceil\frac{N_hn_0}{N_h+n_0-1}\right\rceil.
\end{equation}
Utiliza-se o valor não arredondado de $n_0$ na fórmula. Se $N_h=20$, o resultado para $e=0{,}10$ é 17; se $N_h=100$, é 50; se $N_h=1000$, é 88. Pressupõe-se sorteio sem reposição e população adequadamente coberta. A precisão nominal não corrige vieses de preservação, OCR ou seleção do acervo.

\subsection{Dimensionamento para detecção de uma forma rara}
Uma pergunta mais adequada à busca de atestações é: qual a chance de encontrar pelo menos um documento positivo? Sob ensaios independentes com probabilidade constante $p$, tem-se
\begin{equation}
P(\text{detecção}\mid n,p)=1-(1-p)^n,
\qquad n\geq\left\lceil\frac{\log(1-q)}{\log(1-p)}\right\rceil,
\end{equation}
em que $q$ é a probabilidade de detecção desejada. Para $q=0{,}95$ e $p=0{,}05$, são necessários 59 documentos; para $p=0{,}01$, 299. Com três documentos, a chance é apenas 14,3\% no primeiro cenário e 3,0\% no segundo. Seis documentos elevam esses valores a 26,5\% e 5,9\%. A fórmula explicita por que uma quota pequena não fundamenta afirmações de ausência.

\begin{table}[htbp]\centering\small
\caption{Probabilidade de ao menos uma detecção em cenários hipotéticos independentes. Não são resultados observados.}
\begin{tabular}{rrr}\toprule
Documentos ($n$) & $p=0{,}05$ & $p=0{,}01$ \\\midrule
3 & 14,3\% & 3,0\% \\
6 & 26,5\% & 5,9\% \\
30 & 78,5\% & 26,0\% \\
59 & 95,2\% & 44,7\% \\
97 & 99,3\% & 62,3\% \\
299 & 100,0\% & 95,0\% \\
\bottomrule\end{tabular}
\end{table}

A independência é uma aproximação. Notícias reproduzidas, obras do mesmo autor e sessões com os mesmos participantes podem aumentar a correlação. Se necessário, um planejamento por conglomerados deverá considerar efeito de desenho $D$, com ajuste inicial $n_{\mathrm{aj}}=\lceil D n_0\rceil$, seguido de análise compatível com a amostragem efetiva. Não se adota um valor de $D$ como se fosse estimado pelos dados atuais.

\subsection{Incerteza, zeros e análise temporal}
Para uma amostra aleatória simples binomial, o intervalo de Wilson é uma alternativa ao intervalo normal simétrico, particularmente com contagens pequenas \cite{wilson1927}. Com $\hat p=x/n$, seus limites são
\begin{equation}
\frac{\hat p+z^2/(2n)\ \pm\ z\sqrt{\hat p(1-\hat p)/n+z^2/(4n^2)}}{1+z^2/n}.
\end{equation}
O script inclui a fórmula para reprodução do planejamento. Ela não é aplicada como intervalo populacional à coleta por conveniência deste artigo. Quando não há amostra válida num ano, informa-se ``sem atestação confirmada nesta coleta'', e não uma frequência zero da língua. O denominador precisa ser exposto mesmo quando a contagem de positivos é zero.

Com corpus ampliado, comparações de contagens poderão utilizar modelos com compensação pela extensão textual, efeitos de documento e controles de gênero, país e acervo. Modelos de mudança temporal não devem ser ajustados a uma tabela composta apenas de exemplos selecionados. A forma mais antiga encontrada também é uma estatística dependente da cobertura: um acervo maior pode deslocá-la para trás sem que a história real tenha mudado.

\subsection{Cobertura anual efetiva e plano de expansão}
A rodada de 8 e 9 de novembro de 2026 buscou um testemunho adicional por ano ainda descoberto. Foram incluídos \AnnualAdded{} documentos, cobrindo \AnnualCovered{} dos 110 anos de 1900 a 2009. A seleção dirigida priorizou os Anais do Senado Federal brasileiro, os diários parlamentares portugueses e periódicos brasileiros da Hemeroteca Digital. Entre os novos testemunhos, 87 são registros parlamentares e nove são periódicos. Há 73 registros brasileiros e 23 portugueses. As rubricas editoriais registram o riso numa situação narrada ou transcrita. As grafias miméticas representam o som por escrito. Essas categorias foram mantidas distintas. O OCR serviu à localização em fac-símiles, com conferência da passagem na página. As rubricas eletrônicas foram cotejadas com o contexto na transcrição institucional. O ano segue o documento identificado pelo cabeçalho, edição ou catálogo, conforme anotado no localizador. O protocolo visou comprovar ao menos uma ocorrência por ano e privilegiou séries pesquisáveis. A seleção, portanto, concentra gêneros e instituições. Os dados estão disponíveis na \href{https://github.com/murilloed/historia-riso-escrito-portugues/blob/main/data/annual_coverage.csv}{série anual} e no \href{https://github.com/murilloed/historia-riso-escrito-portugues/blob/main/data/annual_supplement.csv}{inventário dos novos testemunhos}, com URL, forma, nível de evidência e localizador.

O planejamento de 97 documentos por estrato seria muito mais extenso que a quota inicial. Para 110 anos e dois suportes, corresponderia a até 21.340 documentos antes de correções por universo finito, indisponibilidade ou conglomerados. Essa ordem de grandeza é uma consequência da precisão escolhida, não uma exigência universal. Um projeto pode adotar precisão menor, estratos decenais ou objetivo apenas de detecção, desde que justifique a decisão e preserve os limites inferenciais.
